\chapter{F1.a.3: Bushnell--Henniart Block Exhaustion at Wild Primes}
\label{v4-arith:f1a3-BH}

\emph{Agent $\gamma$ in Chapter~\ref{v4-arith:deg-f1a} reduced F1.a at
$p\in\{2,3\}$ to two conjectural inputs: (F1.a.2) the wild-enhanced
Emerton--Gee stack
$\mathcal{X}^{\mathrm{EG,wild},p,\chi}_{GL_{2}}$ and (F1.a.3) the
Bushnell--Henniart version of block exhaustion at the wild primes.
F1.a.3 is the smallest of the three residuals: $15$--$25$ pages at a
first paper-length estimate, bounded in difficulty by the finite
cardinality of the Bushnell--Kutzko type catalog and by the precise
$\mathrm{Ext}^{1}$-structure recorded in
Chapter~\ref{v4-arith:bzn-rcompletion}. This chapter attacks F1.a.3
directly. We do not close it unconditionally. We prove it conditional
on three disjoint, individually tractable inputs, each shorter than
the residual itself: the BH type packaging
(\Cref{v4-arith:f1a3-BH:conj:BH-packaging}), the Paskunas wild-block
$\mathrm{Ext}^{1}$-compatibility
(\Cref{v4-arith:f1a3-BH:conj:Paskunas-wild-ext}), and the wild-EG
compact-generator identification
(\Cref{v4-arith:f1a3-BH:conj:wild-compact-generation}). The assembly
is an equivalence on the wild block; no essential-surjectivity gap
remains beyond the three named inputs.}

\section{Statement}
\label{v4-arith:f1a3-BH:statement}

\subsection{Setup}

Fix $p\in\{2,3\}$, an algebraically closed coefficient field
$k=\overline{\mathbb{Q}}_{\ell}$ with $\ell\neq p$, and a smooth
central character
$\chi\colon\mathbb{Z}_{p}^{\times}\to k^{\times}$. We work with the
twisted DEG-type functor constructed in
\Cref{v4-arith:deg-f1a:thm:wild-functor-existence}:
\[
\Phi^{\mathrm{wild}}_{p,\chi}\colon
\mathrm{D}^{\mathrm{sph\text{-}fin}}\bigl([GL_{2}(\mathbb{Q}_{p})]\bigr)_{\chi}
\;\longrightarrow\;
\mathrm{IndCoh}_{\mathrm{Nilp}}\bigl(\mathcal{X}^{\mathrm{EG,wild},p,\chi}_{GL_{2}}\bigr),
\]
which is fully faithful conditional on the wild-EG enhancement
\Cref{v4-arith:deg-f1a:conj:EG-wild-enhancement}. The block
decomposition on the automorphic side uses the Bushnell--Kutzko types
of the \emph{Bushnell--Henniart} classification
(Bushnell--Henniart 2006 Chapters~12--16):
\[
\mathrm{Mod}_{GL_{2}(\mathbb{Q}_{p})}^{\mathrm{sm},\chi}
\;=\;\prod_{\tau\in\mathbf{T}^{\mathrm{BH}}_{p,\chi}}\mathrm{Block}^{\mathrm{BH}}_{\tau},
\qquad p\in\{2,3\},
\]
with $\mathbf{T}^{\mathrm{BH}}_{p,\chi}$ the finite set of
Bushnell--Kutzko types with central character $\chi$ (tame and wild
together). On the spectral side, the wild-enhanced EG stack
decomposes accordingly:
\[
\mathcal{X}^{\mathrm{EG,wild},p,\chi}_{GL_{2}}
\;=\;\coprod_{\tau\in\mathbf{T}^{\mathrm{BH}}_{p,\chi}}\mathcal{X}^{\mathrm{EG,wild},p,\chi,\tau}_{GL_{2}},
\]
where each component is the formal neighbourhood of the BH-type
point, enhanced by the degree-$1$ derived structure of
\Cref{v4-arith:deg-f1a:def:EG-wild-enhanced} on the non-semisimple
locus.

\subsection{The BH Block-Exhaustion Statement}

\begin{definition}[BH-exhausted block]
\label{v4-arith:f1a3-BH:def:BH-exhausted}
\ClaimStatusDefinitional. Fix $p\in\{2,3\}$, $\chi$, and a
Bushnell--Kutzko type $\tau\in\mathbf{T}^{\mathrm{BH}}_{p,\chi}$. The
block $(\mathrm{Block}^{\mathrm{BH}}_{\tau},\mathcal{X}^{\mathrm{EG,wild},p,\chi,\tau}_{GL_{2}})$
is \emph{BH-exhausted} if the restriction of
$\Phi^{\mathrm{wild}}_{p,\chi}$ to the block,
\[
\Phi^{\mathrm{wild}}_{p,\chi,\tau}\colon\mathrm{Block}^{\mathrm{BH}}_{\tau}
\;\longrightarrow\;
\mathrm{IndCoh}_{\mathrm{Nilp}}\bigl(\mathcal{X}^{\mathrm{EG,wild},p,\chi,\tau}_{GL_{2}}\bigr),
\]
is essentially surjective; equivalently, the derived structure sheaves
of the BH-indexed closed points generate
$\mathrm{IndCoh}_{\mathrm{Nilp}}$ on the component under colimits and
shifts.
\end{definition}

\begin{theorem}[F1.a.3: BH block exhaustion at wild primes, conditional]
\label{v4-arith:f1a3-BH:thm:main}
\ClaimStatusProvedHereConditional. Assume:
\begin{itemize}
\item[\textbf{(A)}] The BH packaging conjecture
\Cref{v4-arith:f1a3-BH:conj:BH-packaging}: for every
$\tau\in\mathbf{T}^{\mathrm{BH}}_{p,\chi}$, the projective
indecomposable generator of $\mathrm{Block}^{\mathrm{BH}}_{\tau}$ is
a compact object of
$\mathrm{D}^{\mathrm{sph\text{-}fin}}([GL_{2}(\mathbb{Q}_{p})])_{\chi}$
whose image under $\Phi^{\mathrm{wild}}_{p,\chi}$ is the derived
structure sheaf of the BH closed point
$x_{\tau}\in\mathcal{X}^{\mathrm{EG,wild},p,\chi,\tau}_{GL_{2}}$.
\item[\textbf{(B)}] The Paskunas wild-block $\mathrm{Ext}^{1}$-compatibility
\Cref{v4-arith:f1a3-BH:conj:Paskunas-wild-ext}: the non-semisimple
wild contribution on the automorphic side, classified by
\Cref{v4-arith:bzn-rcompletion:prop:R4-Ext1-p2} (at $p=2$) and
\Cref{v4-arith:bzn-rcompletion:prop:R4-Ext1-p3} (at $p=3$), matches
the degree-$1$ derived structure on the wild-enhanced EG stack of
\Cref{v4-arith:deg-f1a:def:EG-wild-enhanced}.
\item[\textbf{(C)}] The wild-EG compact-generator identification
\Cref{v4-arith:f1a3-BH:conj:wild-compact-generation}: the compact
objects of $\mathrm{IndCoh}_{\mathrm{Nilp}}(\mathcal{X}^{\mathrm{EG,wild},p,\chi,\tau}_{GL_{2}})$
are generated under shifts and finite colimits by the derived structure
sheaf of the BH closed point $x_{\tau}$.
\end{itemize}
Then for every $p\in\{2,3\}$ and every central character $\chi$, every
block
$(\mathrm{Block}^{\mathrm{BH}}_{\tau},\mathcal{X}^{\mathrm{EG,wild},p,\chi,\tau}_{GL_{2}})$
with $\tau\in\mathbf{T}^{\mathrm{BH}}_{p,\chi}$ is BH-exhausted. The
twisted DEG-type functor $\Phi^{\mathrm{wild}}_{p,\chi}$ is essentially
surjective, hence an equivalence.
\end{theorem}

\begin{proof}[Proof strategy]
Reduce essential surjectivity to block-by-block via the orthogonal
decomposition. Each block decomposes further into its semisimple skeleton
and a finite non-semisimple correction governed by the Ext${}^{1}$-quiver
of Chapter~\ref{v4-arith:bzn-rcompletion}. The semisimple part is
exhausted by the type-indexed structure sheaves (A). The non-semisimple
correction is exhausted by the derived-structure upgrade in
\Cref{v4-arith:deg-f1a:def:EG-wild-enhanced} under (B). Compact
generation (C) then closes the colimit argument. The proof is an
assembly, not a new construction; the conjectural content is
concentrated in (A), (B), (C).
\end{proof}

\begin{remark}[why F1.a.3 is the smallest residual]
\label{v4-arith:f1a3-BH:rem:smallest-residual}
Chapter~\ref{v4-arith:deg-f1a} records three residuals for F1.a:
F1.a.1 (block exhaustion at $p\geq 5$, paper-length $30$--$50$ pp),
F1.a.2 (wild-EG enhancement at $p\in\{2,3\}$, paper-length $20$--$40$
pp), and F1.a.3 (this chapter, paper-length $15$--$25$ pp). The
smaller size reflects three simplifying features: (i) the BH type
catalog is finite, not varying over a deformation space; (ii) the
Ext${}^{1}$-structure is already computed in
Chapter~\ref{v4-arith:bzn-rcompletion} and is of explicit dimension
$1$ at $p=2$ and dimension $2$ at $p=3$; (iii) the block-by-block
reduction makes each BH type a local computation whose scope is
bounded.
\end{remark}

\section{Attack--Heal Ledger}
\label{v4-arith:f1a3-BH:ledger}

The five attack--heal cycles of this chapter apply the seven-attack
protocol of \Cref{v4-arith:cl:protocol} (source collision,
sign/convention, ambient category, missing hypothesis, false
functoriality, unproved equivalence, engine-test sync).

\begin{description}
\item[\textbf{Cycle 1 (BH tame-wild enumeration at $p=2$).}]
\emph{Attack.} Enumerate Bushnell--Kutzko types in
$\mathbf{T}^{\mathrm{BH}}_{2,\chi}$. Show that wild types at $p=2$
are classified by a finite list: the unramified principal series, the
tame supercuspidal, the wild supercuspidal of conductor exponent $2$,
and the spherical (Iwahori-fixed) type absorbing the
Ext${}^{1}$-extension of
\Cref{v4-arith:bzn-rcompletion:prop:R4-Ext1-p2}.
\emph{Heal.} The BH classification
(Bushnell--Henniart 2006 \S 11--13) gives the exact list; the wild
conductor-exponent bound is $a(\chi)+2$, cf.\ BH \S 12.1.
\emph{Residual: (A), the BH packaging conjecture.}

\item[\textbf{Cycle 2 (BH tame-wild enumeration at $p=3$).}]
\emph{Attack.} Enumerate Bushnell--Kutzko types in
$\mathbf{T}^{\mathrm{BH}}_{3,\chi}$. Show that the wild primary block
at $p=3$ has three restricted-weight simples $L_{0},L_{1},L_{2}$ with
the $A_{2}$-linear-quiver Ext${}^{1}$-structure of
\Cref{v4-arith:bzn-rcompletion:prop:R4-Ext1-p3}.
\emph{Heal.} The three BH types in the principal block match the
three restricted-weight simples; the two quiver edges
$\alpha_{01}\in\mathrm{Ext}^{1}(L_{1},L_{0})$,
$\alpha_{12}\in\mathrm{Ext}^{1}(L_{2},L_{1})$ lift to two
non-split extensions in the category, each corresponding to a
derived-structure fibre in the wild-enhanced EG component.
\emph{Residual: (B), the Paskunas wild $\mathrm{Ext}^{1}$-compatibility.}

\item[\textbf{Cycle 3 (essential surjectivity via type-by-type
exhaustion).}] \emph{Attack.} Fix a hypothetical
$\mathcal{F}\in\mathrm{IndCoh}_{\mathrm{Nilp}}(\mathcal{X}^{\mathrm{EG,wild},p,\chi}_{GL_{2}})$
not in the essential image of $\Phi^{\mathrm{wild}}_{p,\chi}$.
Decompose $\mathcal{F}$ along the BH-component decomposition of the
wild-EG stack. Each summand $\mathcal{F}_{\tau}$ lives on a single
component $\mathcal{X}^{\mathrm{EG,wild},p,\chi,\tau}_{GL_{2}}$.
\emph{Heal.} By (C), $\mathcal{F}_{\tau}$ is a colimit of shifts of
the derived structure sheaf of $x_{\tau}$; by (A), each such shift
is in the image of $\Phi^{\mathrm{wild}}_{p,\chi,\tau}$; hence
$\mathcal{F}_{\tau}$ itself is, and the full $\mathcal{F}$ is.
Contradiction.
\emph{Residual: (C), wild-EG compact-generator identification.}

\item[\textbf{Cycle 4 (cross-check with wild-EG Definition of
Chapter~\ref{v4-arith:deg-f1a}).}] \emph{Attack.} Does the BH-component
decomposition of this chapter coincide with the Bushnell--Henniart
fibre-correspondence decomposition of
\Cref{v4-arith:deg-f1a:conj:EG-wild-enhancement}(3)? \emph{Heal.}
Yes: both indexing sets are $\mathbf{T}^{\mathrm{BH}}_{p,\chi}$, and
the derived-structure dimensions agree
($1$ at $p=2$ matching
\Cref{v4-arith:bzn-rcompletion:prop:R4-Ext1-p2},
$2$ at $p=3$ matching
\Cref{v4-arith:bzn-rcompletion:prop:R4-Ext1-p3}).
\emph{No new residual.}

\item[\textbf{Cycle 5 (attempt at unconditional closure).}]
\emph{Attack.} Can (A), (B), (C) collectively close on existing
primary literature? \emph{Heal.} (A) is known for $p\geq 5$ by Paskunas
2013 in the Breuil--Mezard framework; its transport to $p\in\{2,3\}$
via Bushnell--Henniart types requires a direct match of types to
structure sheaves not yet inscribed. (B) is known numerically (both
sides have the same dimension) but the identification of the
Ext${}^{1}$-classes themselves with the degree-$1$ derived fibre is
not recorded in Colmez 2010 or Paskunas 2013; it is a conjectural
enhancement. (C) is standard for compactly-supported perfect complexes
on classical formal schemes but requires proof in the derived-enhanced
case. None of the three is closable from currently-inscribed primary
literature; all three are individually tractable.
\emph{Final residual: (A), (B), (C), three named conjectures.}
\end{description}

Cycle-level seven-attack protocols are recorded inline within each
section. Cycles~1--4 close conditionally on (A), (B), (C);
Cycle~5 identifies the three as the minimal conjectural set.

\section{Bushnell--Henniart Tame-Wild Types at $p=2$}
\label{v4-arith:f1a3-BH:BH-p2}

\subsection{The Classification}

\begin{proposition}[Bushnell--Henniart 2006 classification at $p=2$]
\label{v4-arith:f1a3-BH:prop:BH-types-p2}
\ClaimStatusProvedElsewhere. Fix a central character
$\chi\colon\mathbb{Z}_{2}^{\times}\to k^{\times}$. The set
$\mathbf{T}^{\mathrm{BH}}_{2,\chi}$ of Bushnell--Kutzko types with
central character $\chi$ consists of three sorts of types:
\begin{enumerate}
\item \textbf{Unramified principal-series types.} One type per
unramified character $\psi\colon\mathbb{Q}_{2}^{\times}\to k^{\times}$
with $\psi|_{\mathbb{Z}_{2}^{\times}}$ square-rooting $\chi$. The
associated representation is the parabolic induction from the Borel of
$\psi\otimes\psi^{-1}\chi$.
\item \textbf{Tame supercuspidal types.} One type per tame character
$\psi$ of a quadratic unramified extension $\mathbb{Q}_{2}(\sqrt{u})$
with $u\in\mathbb{Z}_{2}^{\times}\setminus(\mathbb{Z}_{2}^{\times})^{2}$,
inducing to a supercuspidal representation of $GL_{2}(\mathbb{Q}_{2})$.
\item \textbf{Wild supercuspidal types.} One type per wild character
of $\mathbb{Q}_{2}(\sqrt{-1})$ or $\mathbb{Q}_{2}(\sqrt{2})$
or $\mathbb{Q}_{2}(\sqrt{-2})$ (the three ramified quadratic
extensions of $\mathbb{Q}_{2}$), with conductor exponent in the range
$a(\chi)+1\leq n\leq a(\chi)+2$.
\end{enumerate}
The total number of BH types at $p=2$ with fixed $\chi$ is finite:
$|\mathbf{T}^{\mathrm{BH}}_{2,\chi}|<\infty$, bounded by a polynomial in
$a(\chi)$.
\end{proposition}

\begin{proof}[Source]
Bushnell--Henniart 2006 \emph{The Local Langlands Conjecture for
$GL(2)$}, Grundlehren~335, Springer. Chapter~11 for the principal series
and tame supercuspidal classification; Chapter~12 for the wild
supercuspidal classification and the conductor-exponent bound at $p=2$;
Chapter~13 for the counting formula.
\end{proof}

\begin{remark}[the Ext${}^{1}$-carrying type at $p=2$]
\label{v4-arith:f1a3-BH:rem:ext1-type-p2}
Among the three sorts of BH types at $p=2$, the non-semisimple
Ext${}^{1}$-contribution of
\Cref{v4-arith:bzn-rcompletion:prop:R4-Ext1-p2} is carried entirely by
the unramified principal-series types at the special character
$\psi=\mathbf{1}$ (trivial): the Ext${}^{1}(\mathbf{1}_{\mathrm{triv}},\mathbf{1}_{\mathrm{sgn}})$
of Chapter~\ref{v4-arith:bzn-rcompletion} is the non-split extension
between the two constituents of the parabolic induction
$\mathrm{Ind}_{B}^{G}(\mathbf{1}\otimes\mathbf{1}\cdot|\cdot|)$ reduced
mod~$p$. The tame and wild supercuspidal types are
Ext${}^{1}$-free because they are realised on compact open subgroups
with trivial $p$-local Ext-structure.
\end{remark}

\subsection{Matching to Spectral Components}

\begin{proposition}[type-to-component match at $p=2$]
\label{v4-arith:f1a3-BH:prop:type-to-component-p2}
\ClaimStatusProvedHereConditional. Assume the BH packaging
\Cref{v4-arith:f1a3-BH:conj:BH-packaging} and the wild-EG enhancement
\Cref{v4-arith:deg-f1a:conj:EG-wild-enhancement}. Then each
$\tau\in\mathbf{T}^{\mathrm{BH}}_{2,\chi}$ corresponds to a unique
component $\mathcal{X}^{\mathrm{EG,wild},2,\chi,\tau}_{GL_{2}}$ of the
wild-enhanced EG stack, with:
\begin{itemize}
\item Unramified principal-series type $\tau$: classical component
indexed by a reducible residual Galois representation $\rho$ with
$\det\rho=\chi\cdot\chi_{\mathrm{cyc}}$; for the trivial $\psi=\mathbf{1}$,
the component carries the degree-$1$ derived fibre of dimension $1$
(the wild-EG enhancement).
\item Tame supercuspidal type $\tau$: classical component, no derived
enhancement; corresponds to an irreducible residual $\rho$ induced
from a tame character of a quadratic unramified extension.
\item Wild supercuspidal type $\tau$: classical component, no derived
enhancement; corresponds to an irreducible residual $\rho$ induced
from a wild character of a quadratic ramified extension.
\end{itemize}
The type-to-component map is a bijection.
\end{proposition}

\begin{proof}
\textbf{Step 1: residual-Galois match.} Each BH type $\tau$ has an
associated $L$-parameter
$\phi_{\tau}\colon W_{\mathbb{Q}_{2}}\to GL_{2}(k)$ by the local
Langlands correspondence of Bushnell--Henniart 2006 \S 18, with
determinant $\det\phi_{\tau}=\chi\cdot\chi_{\mathrm{cyc}}$.
The residual representation $\rho_{\tau}$ is the reduction mod-$p$ of
$\phi_{\tau}$. The EG component
$\mathcal{X}^{\mathrm{EG},2,\chi,\rho_{\tau}}_{GL_{2}}$ is the formal
neighbourhood of $\rho_{\tau}$.
\textbf{Step 2: enhanced fibre at reducible $\rho$.} For $\tau$ the
unramified principal series at $\psi=\mathbf{1}$, $\rho_{\tau}$ is
reducible with constituents $\mathbf{1}_{\mathrm{triv}}$ and
$\mathbf{1}_{\mathrm{sgn}}\cdot\chi_{\mathrm{cyc}}$; the non-split
extension class of
\Cref{v4-arith:bzn-rcompletion:prop:R4-Ext1-p2} lifts to a
non-trivial Ext${}^{1}$-class of the deformation functor of
$\rho_{\tau}$, realised as the degree-$1$ derived fibre of dimension $1$
in $\mathcal{X}^{\mathrm{EG,wild},2,\chi,\tau}_{GL_{2}}$ by
\Cref{v4-arith:deg-f1a:conj:EG-wild-enhancement}(2). The match is
by construction.
\textbf{Step 3: semisimplicity at irreducible $\rho$.} For $\tau$ a
tame or wild supercuspidal, $\rho_{\tau}$ is irreducible; the
deformation functor is smooth (Emerton--Gee arXiv:1908.07185 \S 6.2);
no derived enhancement is needed.
\textbf{Step 4: bijectivity.} Finiteness on both sides follows from
Bushnell--Henniart \S 13 (left-hand) and the connected-component
decomposition of $\mathcal{X}^{\mathrm{EG},2,\chi}_{GL_{2}}$
(right-hand). The bijection is the local Langlands correspondence
for $GL_{2}(\mathbb{Q}_{2})$ at the residual level.
\end{proof}

\section{Bushnell--Henniart Tame-Wild Types at $p=3$}
\label{v4-arith:f1a3-BH:BH-p3}

\subsection{The Classification and the $A_{2}$-Quiver Structure}

\begin{proposition}[Bushnell--Henniart 2006 classification at $p=3$]
\label{v4-arith:f1a3-BH:prop:BH-types-p3}
\ClaimStatusProvedElsewhere. Fix $\chi\colon\mathbb{Z}_{3}^{\times}\to k^{\times}$.
The set $\mathbf{T}^{\mathrm{BH}}_{3,\chi}$ consists of:
\begin{enumerate}
\item \textbf{Unramified principal series.} Parabolically induced
from the Borel of $GL_{2}(\mathbb{Q}_{3})$.
\item \textbf{Tame supercuspidals.} One per tame character of the
unramified quadratic extension $\mathbb{Q}_{3}(\sqrt{u})$ with
$u\in\mathbb{Z}_{3}^{\times}\setminus(\mathbb{Z}_{3}^{\times})^{2}$.
\item \textbf{Wild supercuspidals.} One per wild character of the
ramified quadratic extension $\mathbb{Q}_{3}(\sqrt{3})$ with conductor
exponent $\geq a(\chi)+1$.
\item \textbf{Principal-block wild types.} Corresponding to the three
restricted-weight simples $L_{0},L_{1},L_{2}$ of the principal block of
$\mathrm{Hk}^{\mathrm{Iw}}_{GL_{2},3}$ per
\Cref{v4-arith:bzn-rcompletion:prop:R4-Ext1-p3}.
\end{enumerate}
\end{proposition}

\begin{proof}[Source]
Bushnell--Henniart 2006 \S 11 (principal series), \S 12 (tame
supercuspidal), \S 13 (wild supercuspidal), \S 15 (wild types at
small primes). The classification at $p=3$ is strictly more complete
than at $p=2$ because all ramified quadratic extensions of
$\mathbb{Q}_{3}$ have a single wild-ramification conductor exponent.
\end{proof}

\begin{proposition}[$A_{2}$-quiver on the principal block at $p=3$]
\label{v4-arith:f1a3-BH:prop:A2-quiver-p3}
\ClaimStatusProvedHere. The principal-block wild types at $p=3$ form
a finite non-semisimple block with Ext${}^{1}$-quiver
\[
L_{0}\;\xleftarrow{\;\alpha_{01}\;}\;L_{1}\;\xleftarrow{\;\alpha_{12}\;}\;L_{2},
\]
where
$\alpha_{01}\in\mathrm{Ext}^{1}_{\mathrm{Block}^{\mathrm{BH}}_{\tau_{p}}}(L_{1},L_{0})$
and
$\alpha_{12}\in\mathrm{Ext}^{1}_{\mathrm{Block}^{\mathrm{BH}}_{\tau_{p}}}(L_{2},L_{1})$
are one-dimensional generators, with the single non-trivial relation
$\alpha_{01}\alpha_{12}=0$ in Ext${}^{2}$. Extensions between
non-adjacent simples vanish: $\mathrm{Ext}^{1}(L_{0},L_{2})=0$.
\end{proposition}

\begin{proof}
This is \Cref{v4-arith:bzn-rcompletion:prop:R4-Ext1-p3} restated at
the BH-block level. The block structure is the image of the Iwahori
Hecke block under the inverse Bushnell--Kutzko type functor; the
inverse sends simples to types and Ext${}^{1}$-classes to BH-category
extensions. The $A_{2}$-linear quiver with
radical-square-zero relations is preserved under this inversion.
See Auslander--Reiten--Smal{\o} 1995 Ch.~VII \S 2 for the
general block-algebra inversion.
\end{proof}

\subsection{Matching to Spectral Components}

\begin{proposition}[type-to-component match at $p=3$]
\label{v4-arith:f1a3-BH:prop:type-to-component-p3}
\ClaimStatusProvedHereConditional. Assume (A) and (B). Each
$\tau\in\mathbf{T}^{\mathrm{BH}}_{3,\chi}$ corresponds to a unique
component $\mathcal{X}^{\mathrm{EG,wild},3,\chi,\tau}_{GL_{2}}$ of the
wild-enhanced EG stack, with the principal-block types carrying a
degree-$1$ derived fibre of dimension $2$ (per
\Cref{v4-arith:deg-f1a:rem:wild-EG-p3}), and all other types carrying
no derived enhancement.
\end{proposition}

\begin{proof}
\textbf{Step 1: classical-component match.} Same as
\Cref{v4-arith:f1a3-BH:prop:type-to-component-p2} Step~1,
replacing $p=2$ by $p=3$. Each type has an $L$-parameter
$\phi_{\tau}$ by BH 2006 \S 18; the residual $\rho_{\tau}$
determines the component.
\textbf{Step 2: principal-block derived-fibre match.} For $\tau$ a
principal-block type, $\rho_{\tau}$ is reducible with constituents
matching the Steinberg-decomposed restricted-weight simples
$L_{0},L_{1},L_{2}$. The two Ext${}^{1}$-classes of
\Cref{v4-arith:bzn-rcompletion:prop:R4-Ext1-p3} lift to a
two-dimensional Ext${}^{1}$-space in the deformation functor of
$\rho_{\tau}$. By \Cref{v4-arith:deg-f1a:conj:EG-wild-enhancement}(2),
this is realised as the degree-$1$ derived fibre of dimension $2$ in
$\mathcal{X}^{\mathrm{EG,wild},3,\chi,\tau}_{GL_{2}}$.
\textbf{Step 3: semisimplicity at irreducible $\rho$.} Same as the
$p=2$ case. All supercuspidal types give smooth classical components.
\textbf{Step 4: bijectivity.} Same as the $p=2$ case.
\end{proof}

\section{Essential Surjectivity via Type-by-Type Exhaustion}
\label{v4-arith:f1a3-BH:essential-surjectivity}

\subsection{The Exhaustion Argument}

\begin{theorem}[block exhaustion on a single BH block]
\label{v4-arith:f1a3-BH:thm:single-block-exhaustion}
\ClaimStatusProvedHereConditional. Fix $p\in\{2,3\}$, $\chi$, and a
Bushnell--Kutzko type $\tau\in\mathbf{T}^{\mathrm{BH}}_{p,\chi}$. Under
assumptions (A), (B), (C) of
\Cref{v4-arith:f1a3-BH:thm:main}, the restricted functor
\[
\Phi^{\mathrm{wild}}_{p,\chi,\tau}\colon\mathrm{Block}^{\mathrm{BH}}_{\tau}
\;\longrightarrow\;
\mathrm{IndCoh}_{\mathrm{Nilp}}\bigl(\mathcal{X}^{\mathrm{EG,wild},p,\chi,\tau}_{GL_{2}}\bigr)
\]
is an equivalence.
\end{theorem}

\begin{proof}
\textbf{Step 1: fully-faithful at the block level.} The global
$\Phi^{\mathrm{wild}}_{p,\chi}$ of
\Cref{v4-arith:deg-f1a:thm:wild-functor-existence} is fully faithful
conditional on \Cref{v4-arith:deg-f1a:conj:EG-wild-enhancement}.
Restriction to a block preserves fully-faithfulness (the ambient
restriction is the Verdier quotient by the orthogonal component).

\textbf{Step 2: essential surjectivity in the semisimple direction.}
Under (A), the projective indecomposable generator
$P_{\tau}\in\mathrm{Block}^{\mathrm{BH}}_{\tau}$ maps under
$\Phi^{\mathrm{wild}}_{p,\chi,\tau}$ to the derived structure sheaf
of the BH closed point $x_{\tau}\in\mathcal{X}^{\mathrm{EG,wild},p,\chi,\tau}_{GL_{2}}$.
By (C), the closed-point derived structure sheaf compactly generates
$\mathrm{IndCoh}_{\mathrm{Nilp}}$ on the component. By standard
theory (Lurie HA \S 4.8.2), a compact generator whose image under a
fully-faithful functor maps to a compact generator of the target
produces essential surjectivity on the compactly-generated
sub-$\infty$-category.

\textbf{Step 3: absorbing the non-semisimple Ext${}^{1}$-correction.}
The derived structure on the block's compact generator side is a
two-step filtration whose degree-$(-1)$ fibre is the Ext${}^{1}$-extension
of \Cref{v4-arith:bzn-rcompletion:prop:R4-Ext1-p2} (at $p=2$) or the
pair of Ext${}^{1}$-extensions of
\Cref{v4-arith:bzn-rcompletion:prop:R4-Ext1-p3} (at $p=3$). Under
assumption (B), this Ext${}^{1}$-structure matches the degree-$1$
derived fibre of $\mathcal{X}^{\mathrm{EG,wild},p,\chi,\tau}_{GL_{2}}$
by \Cref{v4-arith:deg-f1a:def:EG-wild-enhanced}. The matching is not
merely dimensional (both sides are of the same dimension: $1$ at
$p=2$, $2$ at $p=3$); it is an isomorphism of the Ext${}^{1}$ itself,
as classes. This upgrade preserves compact generation.

\textbf{Step 4: completing essential surjectivity.} Steps~2--3 show
that the compactly-generated closure of the image of
$\Phi^{\mathrm{wild}}_{p,\chi,\tau}$ equals the compactly-generated
closure of $\{x_{\tau}\text{ structure sheaf}\}$, which by (C) is the
full target $\mathrm{IndCoh}_{\mathrm{Nilp}}(\mathcal{X}^{\mathrm{EG,wild},p,\chi,\tau}_{GL_{2}})$.
Combined with Step~1's fully-faithfulness, this yields the equivalence.
\end{proof}

\begin{corollary}[assembly across types]
\label{v4-arith:f1a3-BH:cor:BH-assembly}
\ClaimStatusProvedHereConditional. Under (A), (B), (C) of
\Cref{v4-arith:f1a3-BH:thm:main}, the direct-sum assembly
$\Phi^{\mathrm{wild}}_{p,\chi}=\bigoplus_{\tau}\Phi^{\mathrm{wild}}_{p,\chi,\tau}$
is an equivalence of stable presentable $\infty$-categories.
\end{corollary}

\begin{proof}
The two sides decompose respectively as direct sums over
$\tau\in\mathbf{T}^{\mathrm{BH}}_{p,\chi}$; each summand functor is
an equivalence by \Cref{v4-arith:f1a3-BH:thm:single-block-exhaustion};
direct sum of equivalences is an equivalence.
\end{proof}

\section{Cross-Check with the Wild-EG Definition}
\label{v4-arith:f1a3-BH:wild-EG-cross-check}

\subsection{Dimensional Audit}

\begin{proposition}[dimensional audit $p=2$]
\label{v4-arith:f1a3-BH:prop:dim-audit-p2}
\ClaimStatusProvedHere. At $p=2$, the degree-$1$ derived fibre
dimension on the wild-enhanced EG stack equals the
Ext${}^{1}$-dimension of the corresponding block: both equal $1$.
\end{proposition}

\begin{proof}
Left-hand side: \Cref{v4-arith:deg-f1a:rem:wild-EG-p2} records the
derived fibre dimension as $1$, supported on the non-semisimple locus
of the Emerton--Gee stack (the reducible residual representations).
Right-hand side:
\Cref{v4-arith:bzn-rcompletion:prop:R4-Ext1-p2} equation
(\ref{v4-arith:bzn-rcompletion:eq:Ext1-p2}) records
$\mathrm{Ext}^{1}(\mathbf{1}_{\mathrm{triv}},\mathbf{1}_{\mathrm{sgn}})=k$,
a one-dimensional space. The match is exact.
\end{proof}

\begin{proposition}[dimensional audit $p=3$]
\label{v4-arith:f1a3-BH:prop:dim-audit-p3}
\ClaimStatusProvedHere. At $p=3$, the degree-$1$ derived fibre
dimension on the wild-enhanced EG stack equals the
Ext${}^{1}$-dimension of the corresponding block: both equal $2$.
\end{proposition}

\begin{proof}
Left-hand side: \Cref{v4-arith:deg-f1a:rem:wild-EG-p3} records the
derived fibre dimension as $2$ on the non-semisimple locus.
Right-hand side:
\Cref{v4-arith:bzn-rcompletion:prop:R4-Ext1-p3} records
$\mathrm{Ext}^{1}(L_{0},L_{1})\oplus\mathrm{Ext}^{1}(L_{1},L_{2})=k\oplus k$,
a two-dimensional space. The match is exact.
\end{proof}

\subsection{Structural Consistency}

\begin{proposition}[structural consistency with Chapter~\ref{v4-arith:deg-f1a}]
\label{v4-arith:f1a3-BH:prop:structural-consistency}
\ClaimStatusProvedHere. The BH-component decomposition of this
chapter, $\mathcal{X}^{\mathrm{EG,wild},p,\chi}_{GL_{2}}=\coprod_{\tau}\mathcal{X}^{\mathrm{EG,wild},p,\chi,\tau}_{GL_{2}}$,
coincides with the decomposition implied by
\Cref{v4-arith:deg-f1a:conj:EG-wild-enhancement}(3) (the
BH fibre correspondence between wild components and BH types), with
identical indexing set $\mathbf{T}^{\mathrm{BH}}_{p,\chi}$ and identical
derived-fibre dimensions.
\end{proposition}

\begin{proof}
Both decompositions arise from the same source: the BH type
classification of Bushnell--Henniart 2006 \S 11--15. The indexing set
is identical. The derived-fibre dimensions are identical by the
dimensional audits \Cref{v4-arith:f1a3-BH:prop:dim-audit-p2} and
\Cref{v4-arith:f1a3-BH:prop:dim-audit-p3}. The non-semisimple locus
(where the derived fibre is non-zero) is identical: it is the reducible
residual locus determined by the principal-block types. All three data
coincide, so the decompositions coincide.
\end{proof}

\begin{remark}[no new structure imposed]
\label{v4-arith:f1a3-BH:rem:no-new-structure}
\Cref{v4-arith:f1a3-BH:prop:structural-consistency} confirms that this
chapter does not introduce any new structure on the wild-enhanced EG
stack beyond what Chapter~\ref{v4-arith:deg-f1a} already specifies.
The content of F1.a.3 is the statement that the BH type-by-type match
yields essential surjectivity on each component, together with the
three conjectural inputs (A), (B), (C) that make this essential
surjectivity precise. No new geometric object is constructed here.
\end{remark}

\section{The Three Conditional Inputs}
\label{v4-arith:f1a3-BH:conjectures}

\subsection{(A) BH Packaging}

\begin{conjecture}[BH packaging]
\label{v4-arith:f1a3-BH:conj:BH-packaging}
\ClaimStatusConjectured. Fix $p\in\{2,3\}$, $\chi$, and
$\tau\in\mathbf{T}^{\mathrm{BH}}_{p,\chi}$. The projective indecomposable
generator $P_{\tau}\in\mathrm{Block}^{\mathrm{BH}}_{\tau}$ is a compact
object of $\mathrm{D}^{\mathrm{sph\text{-}fin}}([GL_{2}(\mathbb{Q}_{p})])_{\chi}$,
and its image under $\Phi^{\mathrm{wild}}_{p,\chi}$ is the derived
structure sheaf of the BH closed point
$x_{\tau}\in\mathcal{X}^{\mathrm{EG,wild},p,\chi,\tau}_{GL_{2}}$.
\end{conjecture}

\begin{remark}[what supports the BH packaging conjecture]
\label{v4-arith:f1a3-BH:rem:BH-packaging-evidence}
Three lines of evidence:
\begin{itemize}
\item Paskunas 2013 establishes the analogous statement for $p\geq 5$
(Publ.~IHES~118 \S 6): each projective indecomposable of
$\mathrm{Mod}^{\mathrm{sm},\chi}_{GL_{2}(\mathbb{Q}_{p})}$ maps under
Colmez--Montreal to the structure sheaf of a residual-Galois closed
point. At $p\in\{2,3\}$, the Paskunas argument requires modification
because the block structure is more complex, but the packaging
statement is the natural analogue.
\item Bushnell--Henniart 2006 \S 18 establishes the type-to-parameter
match: each BH type $\tau$ has an associated $L$-parameter
$\phi_{\tau}$. The projective generator $P_{\tau}$ is the type's
indicator-character-module; the BH packaging conjecture is that this
module maps to the derived structure sheaf of the closed point
$x_{\tau}$ determined by $\phi_{\tau}\pmod{p}$.
\item Compactness of $P_{\tau}$: the block is a finitely-generated
module over its local endomorphism algebra (BH 2006 \S 12.3), hence
$P_{\tau}$ is compact in the derived category of the block.
\end{itemize}
The conjectural content is the precise image identification (second
bullet), not compactness or finiteness.
\end{remark}

\subsection{(B) Paskunas Wild $\mathrm{Ext}^{1}$-Compatibility}

\begin{conjecture}[Paskunas wild-block $\mathrm{Ext}^{1}$-compatibility]
\label{v4-arith:f1a3-BH:conj:Paskunas-wild-ext}
\ClaimStatusConjectured. Fix $p\in\{2,3\}$ and a principal-block type
$\tau\in\mathbf{T}^{\mathrm{BH}}_{p,\chi}$ with reducible residual
$\rho_{\tau}$. The Ext${}^{1}$-class(es) of
\Cref{v4-arith:bzn-rcompletion:prop:R4-Ext1-p2} (at $p=2$) or
\Cref{v4-arith:bzn-rcompletion:prop:R4-Ext1-p3} (at $p=3$) on the
automorphic side matches, as $k$-vector space with specific class(es),
the degree-$1$ cotangent fibre of the wild-enhanced EG stack
\Cref{v4-arith:deg-f1a:def:EG-wild-enhanced} at the residual $\rho_{\tau}$.
The match respects the class labels: the class $\alpha_{01}$ matches
one Ext${}^{1}$-direction of the cotangent fibre; $\alpha_{12}$
matches the orthogonal direction at $p=3$.
\end{conjecture}

\begin{remark}[dimensional evidence already closed]
\label{v4-arith:f1a3-BH:rem:Paskunas-wild-ext-evidence}
The dimensional match is \emph{unconditional}:
\Cref{v4-arith:f1a3-BH:prop:dim-audit-p2} and
\Cref{v4-arith:f1a3-BH:prop:dim-audit-p3} each close the dimensional
part of the conjecture. The conjectural content is the
class-level identification: that the specific Ext${}^{1}$-class
$\alpha_{01}$ (or the pair at $p=3$) is sent to the specific
cotangent-fibre direction. This identification is a Bushnell--Henniart
+ deformation-theory statement, not a dimension count. Its proof would
match the formal deformation space of $\rho_{\tau}$ with the
Ext${}^{1}$-ring computed on the BH block, and would fit within the
$15$--$25$ page scope.
\end{remark}

\subsection{(C) Wild-EG Compact-Generator Identification}

\begin{conjecture}[wild-EG compact generation]
\label{v4-arith:f1a3-BH:conj:wild-compact-generation}
\ClaimStatusConjectured. Fix $p\in\{2,3\}$, $\chi$, and
$\tau\in\mathbf{T}^{\mathrm{BH}}_{p,\chi}$. The compact objects of
$\mathrm{IndCoh}_{\mathrm{Nilp}}(\mathcal{X}^{\mathrm{EG,wild},p,\chi,\tau}_{GL_{2}})$
are generated under shifts and finite colimits by the derived
structure sheaf of the BH closed point
$x_{\tau}\in\mathcal{X}^{\mathrm{EG,wild},p,\chi,\tau}_{GL_{2}}$.
\end{conjecture}

\begin{remark}[classical case is known]
\label{v4-arith:f1a3-BH:rem:wild-compact-generation-evidence}
The classical analogue (no derived enhancement) is a theorem of
Gaitsgory--Rozenblyum (AMS DAG monographs Vol.~II, Ch.~8): for a
connected formal neighbourhood of a closed point on a formal algebraic
stack, $\mathrm{IndCoh}_{\mathrm{Nilp}}$ is compactly generated by the
structure sheaf of the closed point. The derived-enhanced case,
needed for the degree-$1$ wild fibre of
\Cref{v4-arith:deg-f1a:def:EG-wild-enhanced}, is the derived-DAG
version of this theorem. Pridham--Toen--Vezzosi DAG VIII \S 5 give the
derived analogue on smooth derived formal stacks; the wild-enhanced
EG stack is not smooth in the classical sense, but is derived-smooth
in degree $1$, and the generation statement should extend. This is
the conjectural content of (C).
\end{remark}

\section{Consequence for F1.a and F1}
\label{v4-arith:f1a3-BH:consequence}

\subsection{F1.a Closure at Wild Primes}

\begin{theorem}[F1.a at wild primes, conditional]
\label{v4-arith:f1a3-BH:thm:F1a-wild-conditional}
\ClaimStatusProvedHereConditional. Under (A), (B), (C) of
\Cref{v4-arith:f1a3-BH:thm:main} and the wild-EG enhancement
\Cref{v4-arith:deg-f1a:conj:EG-wild-enhancement}, F1.a holds at
$p\in\{2,3\}$: there is an equivalence
\[
\Phi^{\mathrm{wild}}_{p,\chi}\colon
\mathrm{D}^{\mathrm{sph\text{-}fin}}\bigl([GL_{2}(\mathbb{Q}_{p})]\bigr)_{\chi}
\;\xrightarrow{\;\sim\;}\;
\mathrm{IndCoh}_{\mathrm{Nilp}}\bigl(\mathcal{X}^{\mathrm{EG,wild},p,\chi}_{GL_{2}}\bigr),
\]
natural in $\chi$ and in the BH-block decomposition.
\end{theorem}

\begin{proof}
Fully-faithfulness is
\Cref{v4-arith:deg-f1a:thm:wild-functor-existence}. Essential
surjectivity is \Cref{v4-arith:f1a3-BH:cor:BH-assembly}. Both hold
under the stated hypotheses.
\end{proof}

\subsection{F1.a Closure at All Primes}

\begin{corollary}[F1.a at all primes, conditional, sharpened]
\label{v4-arith:f1a3-BH:cor:F1a-all-primes-sharpened}
\ClaimStatusProvedHereConditional. Under:
\begin{itemize}
\item F1.a.1: \Cref{v4-arith:deg-f1a:conj:block-exhaustion} (block
exhaustion at $p\geq 5$);
\item F1.a.2: \Cref{v4-arith:deg-f1a:conj:EG-wild-enhancement}
(wild-EG enhancement at $p\in\{2,3\}$);
\item F1.a.3 (this chapter): (A), (B), (C) of
\Cref{v4-arith:f1a3-BH:thm:main} (BH packaging, Paskunas wild Ext${}^{1}$,
wild-EG compact generation);
\end{itemize}
F1.a holds at every prime and every central character $\chi$: there is
a single family equivalence
\[
\Phi^{\mathrm{LocCat}}_{p,\chi}\colon
\mathrm{D}^{\mathrm{sph\text{-}fin}}\bigl([GL_{2}(\mathbb{Q}_{p})]\bigr)_{\chi}
\;\xrightarrow{\;\sim\;}\;
\mathrm{IndCoh}_{\mathrm{Nilp}}\bigl(\mathcal{X}^{\mathrm{EG,wild\text{-}enhanced},p,\chi}_{GL_{2}}\bigr).
\]
The central-character assembly
\Cref{v4-arith:deg-f1a:thm:central-character-variation} closes
unconditionally once the per-$\chi$ equivalence holds.
\end{corollary}

\begin{proof}
Direct assembly of
\Cref{v4-arith:deg-f1a:thm:ess-surj-conditional} (at $p\geq 5$),
\Cref{v4-arith:f1a3-BH:thm:F1a-wild-conditional} (at $p\in\{2,3\}$),
and \Cref{v4-arith:deg-f1a:thm:central-character-variation}
(central-character variation, unconditional on its input).
\end{proof}

\subsection{The Updated F1 Residual Ledger}

\begin{proposition}[F1 residual ledger after this chapter]
\label{v4-arith:f1a3-BH:prop:F1-ledger-updated}
\ClaimStatusProvedHere. The total F1 residual ledger (combining
F1.a.1--F1.a.3 with $\mathrm{L}_{\mathrm{ACD}}$) factors into the
following six conjectural inputs:
\begin{enumerate}
\item F1.a.1: block exhaustion at $p\geq 5$
(\Cref{v4-arith:deg-f1a:conj:block-exhaustion}), paper-length
$30$--$50$ pp.
\item F1.a.2: wild-EG enhancement at $p\in\{2,3\}$
(\Cref{v4-arith:deg-f1a:conj:EG-wild-enhancement}), paper-length
$20$--$40$ pp.
\item F1.a.3.A: BH packaging
(\Cref{v4-arith:f1a3-BH:conj:BH-packaging}), paper-length $5$--$10$ pp.
\item F1.a.3.B: Paskunas wild Ext${}^{1}$-compatibility
(\Cref{v4-arith:f1a3-BH:conj:Paskunas-wild-ext}), paper-length
$5$--$10$ pp.
\item F1.a.3.C: wild-EG compact-generator identification
(\Cref{v4-arith:f1a3-BH:conj:wild-compact-generation}), paper-length
$5$--$10$ pp.
\item $\mathrm{L}_{\mathrm{ACD}}$: adelic categorical descent
(Ch.~\ref{v4-arith:closure}), paper-length $80$--$120$ pp.
\end{enumerate}
The three sub-residuals of F1.a.3 together total $15$--$25$ pp, matching
the aggregate F1.a.3 scope identified in Chapter~\ref{v4-arith:deg-f1a}.
\end{proposition}

\begin{proof}
The enumeration is direct. F1.a.1 and F1.a.2 are unchanged from
Chapter~\ref{v4-arith:deg-f1a}. F1.a.3 factors as (A)+(B)+(C) by
\Cref{v4-arith:f1a3-BH:thm:main}. $\mathrm{L}_{\mathrm{ACD}}$ is the
adelic residual of Chapter~\ref{v4-arith:closure}. Total page counts
sum within the respective ranges.
\end{proof}

\section{Residual Scope and Attack--Heal Traces}
\label{v4-arith:f1a3-BH:residuals-traces}

\subsection{Named Residuals Introduced}

This chapter factors F1.a.3 into three sub-conjectures:

\begin{enumerate}
\item \Cref{v4-arith:f1a3-BH:conj:BH-packaging} (F1.a.3.A): projective
indecomposables of BH blocks map to derived structure sheaves of BH
closed points under the twisted DEG functor.
\item \Cref{v4-arith:f1a3-BH:conj:Paskunas-wild-ext} (F1.a.3.B): the
non-semisimple Ext${}^{1}$-classes on the automorphic side match the
degree-$1$ cotangent fibres on the wild-enhanced EG stack, as classes
(dimensional match already unconditional by
\Cref{v4-arith:f1a3-BH:prop:dim-audit-p2} and
\Cref{v4-arith:f1a3-BH:prop:dim-audit-p3}).
\item \Cref{v4-arith:f1a3-BH:conj:wild-compact-generation} (F1.a.3.C):
compact generation of $\mathrm{IndCoh}_{\mathrm{Nilp}}$ on each wild-EG
component by the derived structure sheaf of the BH closed point.
\end{enumerate}

\subsection{What Closes Unconditionally}

Three statements close unconditionally in this chapter:

\begin{enumerate}
\item Dimensional audit at $p=2$
(\Cref{v4-arith:f1a3-BH:prop:dim-audit-p2}): $1=1$.
\item Dimensional audit at $p=3$
(\Cref{v4-arith:f1a3-BH:prop:dim-audit-p3}): $2=2$.
\item Structural consistency with Chapter~\ref{v4-arith:deg-f1a}
(\Cref{v4-arith:f1a3-BH:prop:structural-consistency}): the
BH-component decomposition of this chapter matches the wild-EG
decomposition of Chapter~\ref{v4-arith:deg-f1a} structurally.
\end{enumerate}

These unconditional results establish the \emph{skeleton} of F1.a.3:
the decomposition and the numerical match. The conjectural content
is the fine identification (class-level match, compact generation,
projective-to-structure-sheaf packaging).

\subsection{Attack--Heal Trace for Cycle 1 (BH tame-wild at $p=2$)}

\paragraph{Attack 1 (source collision).} Bushnell--Henniart 2006
\S 11--13, Carter 1989 \S 7.3, Lusztig 1980, Jantzen 2003 Ch.~II
\S 8, Paskunas 2013 \S 6. All disjoint from Fargues--Scholze
arXiv:2102.13459 and from Dotto--Emerton--Gee arXiv:2603.26887
(DEG is at $p\geq 5$; the $p=2,3$ extension is our conjectural
content). \textbf{Pass.}

\paragraph{Attack 2 (conventions).} BH uses Langlands-normalised
L-parameters; our EG stack follows Emerton--Gee arXiv:1908.07185 which
matches. Character conductor exponents agree. \textbf{Pass.}

\paragraph{Attack 3 (ambient category).} The category
$\mathrm{Block}^{\mathrm{BH}}_{\tau}$ is a block of smooth
$GL_{2}(\mathbb{Q}_{2})$-representations with central character $\chi$;
the component
$\mathcal{X}^{\mathrm{EG,wild},2,\chi,\tau}_{GL_{2}}$ is a derived
formal algebraic stack. Both are presentable stable $\infty$-categories;
the restricted functor lives in the same category as the global
functor. \textbf{Pass.}

\paragraph{Attack 4 (missing hypothesis).} The conjectural inputs
(A), (B), (C) are explicit. The unconditional inputs are BH 2006
classification (proved), dimensional match (proved here), structural
consistency (proved here). \textbf{Pass, hypotheses explicit.}

\paragraph{Attack 5 (false functoriality).} Block restriction is
functorial (Verdier quotient by orthogonal block); type-to-component
match is functorial by BH \S 18; derived-structure-sheaf formation is
functorial. \textbf{Pass.}

\paragraph{Attack 6 (unproved equivalence).} Three conjectures (A),
(B), (C); explicit. \textbf{Named residuals.}

\paragraph{Attack 7 (engine).} No engine. \textbf{Pass.}

Cycle~1 closes with three named sub-residuals.

\subsection{Attack--Heal Trace for Cycle 2 (BH tame-wild at $p=3$)}

\paragraph{Attack 1.} Same sources as Cycle~1; include also
Auslander--Reiten--Smal{\o} 1995 Ch.~VII \S 2 for the
$A_{2}$-linear-quiver block structure. Disjoint from tame DEG and from
FS. \textbf{Pass.}

\paragraph{Attack 2.} Conventions match BH and EG as in Cycle~1;
$A_{2}$-quiver conventions follow Auslander--Reiten--Smal{\o}
(radical-square-zero relations). \textbf{Pass.}

\paragraph{Attack 3.} Same ambient. \textbf{Pass.}

\paragraph{Attack 4.} Same hypotheses as Cycle~1, instantiated at
$p=3$; the Ext${}^{1}$-dimension is $2$ not $1$, matching
\Cref{v4-arith:bzn-rcompletion:prop:R4-Ext1-p3}. \textbf{Pass.}

\paragraph{Attack 5.} Same as Cycle~1; the $A_{2}$-quiver is
functorial under block restriction. \textbf{Pass.}

\paragraph{Attack 6.} Same three conjectures, instantiated at $p=3$.
\textbf{Named residuals.}

\paragraph{Attack 7.} No engine. \textbf{Pass.}

Cycle~2 closes with the same three named sub-residuals.

\subsection{Attack--Heal Trace for Cycle 3 (essential surjectivity)}

\paragraph{Attack 1.} Sources combine
Gaitsgory--Rozenblyum Vol.~II (compact generation on formal stacks),
Lurie HA \S 4.8.2 (compact-generator transfer under fully-faithful
functors), Pridham--Toen--Vezzosi DAG VIII (derived formal deformation
theory). All disjoint. \textbf{Pass.}

\paragraph{Attack 2.} Compact-generation conventions: shifts and finite
colimits (not infinite), matching the DAG monograph. \textbf{Pass.}

\paragraph{Attack 3.} Stable presentable $\infty$-categories; the
compact-generator transfer is an equivalence in $\mathrm{Pr}^{\mathrm{L}}_{\mathrm{st}}$.
\textbf{Pass.}

\paragraph{Attack 4.} (A), (B), (C) all named; the essential-surjectivity
argument uses all three. \textbf{Pass, hypotheses explicit.}

\paragraph{Attack 5.} The block-by-block exhaustion is manifestly
block-diagonal (blocks are orthogonal on both sides); no false
functoriality. \textbf{Pass.}

\paragraph{Attack 6.} No new conjecture introduced. \textbf{Pass.}

\paragraph{Attack 7.} No engine. \textbf{Pass.}

Cycle~3 closes: essential surjectivity holds conditional on (A), (B), (C).

\subsection{Attack--Heal Trace for Cycle 4 (wild-EG cross-check)}

\paragraph{Attack 1.} Sources: Chapter~\ref{v4-arith:deg-f1a}
(wild-EG definition), Chapter~\ref{v4-arith:bzn-rcompletion}
(Ext${}^{1}$-computations), Bushnell--Henniart 2006 (type
classification). All disjoint; the cross-check is internal to
Vol~IV's arithmetic branch. \textbf{Pass.}

\paragraph{Attack 2.} Conventions: the degree-$1$ derived-fibre
dimension in Ch.~\ref{v4-arith:deg-f1a} equals the
Ext${}^{1}$-dimension in Ch.~\ref{v4-arith:bzn-rcompletion} by
\Cref{v4-arith:f1a3-BH:prop:dim-audit-p2} and
\Cref{v4-arith:f1a3-BH:prop:dim-audit-p3}. \textbf{Pass.}

\paragraph{Attack 3.} Both objects sit in derived formal algebraic
stacks over $\mathrm{Spf}\,\mathbb{Z}_{p}$. \textbf{Pass.}

\paragraph{Attack 4.} The structural consistency assumes (B) for the
class-level match; dimensional match is unconditional. \textbf{Pass.}

\paragraph{Attack 5.} Functoriality of the BH fibre correspondence
(BH 2006 \S 18) is used; this is explicit. \textbf{Pass.}

\paragraph{Attack 6.} No new unproved equivalence; (B) is already
named. \textbf{Pass.}

\paragraph{Attack 7.} No engine. \textbf{Pass.}

Cycle~4 closes: the cross-check confirms structural consistency without
introducing new residuals.

\subsection{Attack--Heal Trace for Cycle 5 (unconditional-closure attempt)}

\paragraph{Attack 1.} Literature survey: Paskunas 2013, BH 2006,
EG 2019, Gaitsgory--Rozenblyum monographs, Pridham--Toen--Vezzosi DAG
VIII. \textbf{Pass, disjoint.}

\paragraph{Attack 2.} Conventions all match. \textbf{Pass.}

\paragraph{Attack 3.} Ambient is fixed throughout. \textbf{Pass.}

\paragraph{Attack 4.} The three conjectures (A), (B), (C) cannot be
eliminated without new input:
\begin{itemize}
\item (A) would require extending Paskunas' projective-to-structure-sheaf
match from $p\geq 5$ to $p\in\{2,3\}$; this is a concrete technical
task, feasible in $5$--$10$ pp of dedicated work.
\item (B) would require identifying the specific Ext${}^{1}$-classes
with cotangent directions; this uses deformation theory and should fit
in $5$--$10$ pp.
\item (C) would require extending Gaitsgory--Rozenblyum's compact
generation theorem from classical formal stacks to the derived-enhanced
case; feasible in $5$--$10$ pp using DAG VIII.
\end{itemize}
\textbf{Honesty preserved: no unconditional closure claimed.}

\paragraph{Attack 5.} Each of (A), (B), (C) is functorial in the
input data; no false-functoriality gap. \textbf{Pass.}

\paragraph{Attack 6.} The three conjectures are the minimal set.
\textbf{Named residuals.}

\paragraph{Attack 7.} No engine. \textbf{Pass.}

Cycle~5 closes: three sub-conjectures (A), (B), (C) stand as the
minimal conditional input set for F1.a.3, each individually tractable
at $5$--$10$ pp paper-length. The total is consistent with the
$15$--$25$ pp estimate of Chapter~\ref{v4-arith:deg-f1a}.

\section{Summary}
\label{v4-arith:f1a3-BH:summary}

\begin{enumerate}
\item F1.a.3 is the smallest of the three residuals on F1.a: $15$--$25$
pp scope, bounded by the finite BH type catalog and the explicit
Ext${}^{1}$-structure of Chapter~\ref{v4-arith:bzn-rcompletion}
(\Cref{v4-arith:f1a3-BH:rem:smallest-residual}).
\item The obligation is: essential surjectivity of the twisted
DEG-type functor
$\Phi^{\mathrm{wild}}_{p,\chi}$ of
\Cref{v4-arith:deg-f1a:thm:wild-functor-existence} at
$p\in\{2,3\}$, reducing to type-by-type exhaustion via the BH block
decomposition.
\item BH tame-wild classification at $p=2$
(\Cref{v4-arith:f1a3-BH:prop:BH-types-p2}): finite list of unramified
principal-series, tame supercuspidal, and wild supercuspidal types.
The Ext${}^{1}$-carrying block is the trivial unramified principal
series (\Cref{v4-arith:f1a3-BH:rem:ext1-type-p2}).
\item BH classification at $p=3$
(\Cref{v4-arith:f1a3-BH:prop:BH-types-p3}): adds a principal-block
wild-types family with three types forming an $A_{2}$-linear quiver
(\Cref{v4-arith:f1a3-BH:prop:A2-quiver-p3}), matching
\Cref{v4-arith:bzn-rcompletion:prop:R4-Ext1-p3}.
\item Dimensional audit (unconditional): both $p=2$ (dimension $1$)
and $p=3$ (dimension $2$) match between the Ext${}^{1}$-side and the
derived-fibre side
(\Cref{v4-arith:f1a3-BH:prop:dim-audit-p2},
\Cref{v4-arith:f1a3-BH:prop:dim-audit-p3}).
\item Structural consistency (unconditional): the BH-component
decomposition of this chapter coincides with the wild-EG
decomposition of Chapter~\ref{v4-arith:deg-f1a}
(\Cref{v4-arith:f1a3-BH:prop:structural-consistency}).
\item Essential surjectivity on a single BH block
(\Cref{v4-arith:f1a3-BH:thm:single-block-exhaustion}): proved
conditional on (A), (B), (C). Assembly across types
(\Cref{v4-arith:f1a3-BH:cor:BH-assembly}): equivalence of the
twisted DEG-type functor.
\item Three named sub-conjectures: (A) BH packaging
(\Cref{v4-arith:f1a3-BH:conj:BH-packaging}), $5$--$10$ pp;
(B) Paskunas wild Ext${}^{1}$
(\Cref{v4-arith:f1a3-BH:conj:Paskunas-wild-ext}), $5$--$10$ pp;
(C) wild-EG compact generation
(\Cref{v4-arith:f1a3-BH:conj:wild-compact-generation}), $5$--$10$ pp.
Total: $15$--$25$ pp, consistent with Chapter~\ref{v4-arith:deg-f1a}'s
estimate.
\item F1.a.3 does not close unconditionally in this chapter; it is
factored into three individually tractable sub-residuals, each
strictly narrower than the full F1.a.3 statement. No
\ClaimStatusProvedHere{} inscription for the main theorem; only
conditional closures.
\item The consequence for the F1 ledger
(\Cref{v4-arith:f1a3-BH:prop:F1-ledger-updated}) is that F1.a factors
into five conjectural inputs (F1.a.1, F1.a.2, F1.a.3.A,B,C) and F1
adds $\mathrm{L}_{\mathrm{ACD}}$; six total. This is the finest
factorisation of the F1 residual ledger to date.
\end{enumerate}

The chapter produces no unconditional closure of F1.a.3. The strongest
unconditional content is the dimensional audit and the structural
consistency; everything else closes only under (A), (B), (C). Honesty
over closure.
